\vspace{-1em}
\section{\uppercase{Method}}
\label{sec:method}
\vspace{-1em}
\subsection{Overview}
\vspace{-1em}

A \textbf{knowledge graph} (KG) is a directed multi-relational graph
$\mathcal{G} = (\mathcal{V}, \mathcal{E}, \mathcal{R})$ with
$\mathcal{V}$ a set of entities, $\mathcal{R}$ a set of relation types,
and $\mathcal{E} \subseteq \mathcal{V} \times \mathcal{R} \times \mathcal{V}$
a set of labelled edges.  Each edge $(\mathbf{s}, r, \mathbf{o}) \in \mathcal{E}$
denotes that subject entity $\mathbf{s}$ is connected to object entity $\mathbf{o}$
via relation $r$.  Entities may carry multimodal attributes: visual features
$\mathbf{x}^{\text{vis}}_i \in \mathbb{R}^{D_v}$, such as image embeddings of artworks,
textual features $\mathbf{x}^{\text{txt}}_i \in \mathbb{R}^{D_t}$, such as text
embeddings of entity names, and structural features learned via a KGE model.

Given a query $(\mathbf{s}, r, ?)$, the goal of \textbf{path-aware reasoning} is to
rank a set of candidate objects $\mathcal{O}_{\text{cand}} = \{\mathbf{o}_1, \ldots, \mathbf{o}_K\}$
by scoring the paths that connect $\mathbf{s}$ to each $\mathbf{o}_k$ within $\mathcal{G}$.
Formally, let $\mathcal{P}(\mathbf{s}, \mathbf{o})$ denote the set of paths from
$\mathbf{s}$ to $\mathbf{o}$, where a path
$p = (\mathbf{v}_0, r_1, \mathbf{v}_1, r_2, \mathbf{v}_2, \ldots, r_\ell, \mathbf{v}_\ell)$
is a sequence of alternating entities and relations with $\mathbf{v}_0 = \mathbf{s}$ and
$\mathbf{v}_\ell = \mathbf{o}$.  The number of edges $\ell$ is the path length (in hops).
A scoring function $f: \mathcal{P} \times \mathcal{R} \to \mathbb{R}$ assigns a real-valued
score $s = f(p, r)$ to each path given the query relation $r$.  The object with the
highest scoring path is selected as the answer.

We propose the \textbf{Path-Aware Multimodal Reasoning (PAMR)} framework, a Transformer-based
architecture that encodes paths as interleaved sequences of entity and relation tokens,
each embedded via a frozen vision-language knowledge graph embedding (VL-KGE) model.
PAMR introduces two key innovations: (i)~a paired sampling scheme that, unlike
conventional uniform negative sampling, yields batches with exactly balanced positive
and negative examples per query, and (ii)~a token representation that
concatenates real and imaginary relation embeddings while using only the fused real entity
representation.  These design choices enable the Transformer to learn discriminative path
representations from frozen KGE token embeddings without modifying the underlying knowledge graph
representations.

\vspace{-1em}
\subsection{Embedding Module}
\vspace{-1em}
\label{sec:embedding-module}

To supply pre-encoded relational structure without incurring additional training cost, we
use a frozen pre-trained \textbf{knowledge graph embedding (KGE)} model.  A KGE model
learns entity embeddings
$\mathbf{E} \in \mathbb{R}^{|\mathcal{V}| \times D}$ and relation embeddings
$\mathbf{R} \in \mathbb{R}^{|\mathcal{R}| \times D}$ such that a scoring function
$\phi: \mathcal{V} \times \mathcal{R} \times \mathcal{V} \to \mathbb{R}$
assigns higher values to true triples than to false ones.  We consider four standard
KGE architectures (ComplEx~\cite{trouillon2016complex}, TransE~\cite{bordes2013translating},
DistMult~\cite{yang2015embedding}, RotatE~\cite{sun2019rotate}), each provisioning a
$D$-dimensional real vector per entity and relation.

Entities in a multimodal KG possess heterogeneous modalities.  Let
$\boldsymbol{\Phi}^{\text{st}} \in \mathbb{R}^{|\mathcal{V}| \times D}$ be the trainable
structural embedding table, $\boldsymbol{\Phi}^{\text{vis}}: \mathcal{V} \to \mathbb{R}^{D}$
the frozen visual projector, and $\boldsymbol{\Phi}^{\text{txt}}: \mathcal{V} \to \mathbb{R}^{D}$
the frozen textual projector.  The fused entity representation of entity $i$ is:
\begin{equation}
\mathbf{e}^{\text{fused}}_i = \frac{1}{|\mathcal{M}(i)|}\sum_{m \in \mathcal{M}(i)} \tilde{\boldsymbol{\Phi}}_m(i)
\label{eq:fused-emb}
\end{equation}
Here $\mathcal{M}(i) \subseteq \{\text{st}, \text{vis}, \text{txt}\}$ denotes the set of modalities
available for entity $i$, and $\tilde{\boldsymbol{\Phi}}_m$ denotes the corresponding
projector after optional linear layers for aligning dimensions.  This per-entity averaging
naturally handles modality asymmetry, where different entity types carry different modality
combinations without requiring a fixed global fusion scheme.

Among the four KGE architectures, ComplEx~\cite{trouillon2016complex} and
RotatE~\cite{sun2019rotate} employ complex-valued embeddings that decompose each entity
and relation into real and imaginary components:
$\mathbf{e} = [\mathbf{e}_{\text{re}}; \mathbf{e}_{\text{im}}] \in \mathbb{R}^{2D}$,
$\mathbf{r} = [\mathbf{r}_{\text{re}}; \mathbf{r}_{\text{im}}] \in \mathbb{R}^{2D}$.
TransE and DistMult are purely real-valued models operating in $\mathbb{R}^D$.  The
scoring functions are: ComplEx, $\operatorname{Re}\langle \mathbf{s}, \mathbf{r}, \overline{\mathbf{o}} \rangle$, which computes an elementwise triple dot product across real and imaginary components, TransE,
$m - \|\mathbf{s} + \mathbf{r} - \mathbf{o}\|_1$, DistMult, $\sum_i \mathbf{s}_i \mathbf{r}_i \mathbf{o}_i$,
and RotatE, $m - \|\mathbf{s} \circ \mathbf{r} - \mathbf{o}\|$ with $\mathbf{r}_i = \cos\theta_i + \jmath\sin\theta_i$.

Complex-valued KGE models output embeddings of twice the real-valued dimension.
To decide whether this extra capacity carries useful signal for the Transformer, we
examine the real and imaginary components independently.

\vspace{-1em}
\paragraph{Real and imaginary parts in PAMR.}
The imaginary entity component in ComplEx and RotatE is constructed via a gated linear projection
$\mathbf{e}_{\text{im}} = \tanh(g) \cdot \mathbf{W}_{\text{vi}}(\mathbf{e}_{\text{re}})$.
Here $g \in \mathbb{R}$ is a learned scalar gate and $\mathbf{W}_{\text{vi}} \in \mathbb{R}^{D \times D}$
is a trainable projection.  This construction suffers from an initialization deadlock:
$\mathbf{W}_{\text{vi}}$ is initialised to zero and $g$ to zero, so each parameter's gradient
depends on the other being nonzero, so neither escapes during training.  Consequently, for
entities that rely on this projection, such as entities that lack their own textual embedding table,
the imaginary part remains at its near-zero initialized value.  Concatenating near-zero
vectors would inject noise into the Transformer input.  PAMR therefore uses \textbf{only the
fused real part} for entity tokens: $\mathbf{e}^{\text{raw}}_i = \mathbf{e}^{\text{fused}}_{i,\text{re}}$.

For relations, however, the imaginary embeddings are learned from separate, fully-trained
embedding tables $\mathbf{R}_{\text{im}} \in \mathbb{R}^{|\mathcal{R}| \times D}$ for ComplEx
or via a mapping of angles to complex values $[\cos\boldsymbol{\theta}; \sin\boldsymbol{\theta}]$ for RotatE,
neither of which suffers from the gating deadlock.  PAMR therefore concatenates both components
for relation tokens:
\begin{equation}
\mathbf{r}^{\text{raw}} = \begin{cases}
    [\,\mathbf{r}_{\text{re}}; \mathbf{r}_{\text{im}}\,] \in \mathbb{R}^{2D}
        & \text{(ComplEx, RotatE)} \\
    \mathbf{r}_{\text{re}} \in \mathbb{R}^{D}
        & \text{(TransE, DistMult)}
\end{cases}
\label{eq:rel-raw}
\end{equation}

\vspace{-1em}
\paragraph{Dimension automatic detection.}
Let $d_{\mathcal{V}}$ and $d_{\mathcal{R}}$ denote the effective dimensionality of entity and
relation representations, derived from the KGE dimension $D$:
\begin{equation}
d_{\mathcal{V}} = D,\qquad
d_{\mathcal{R}} = \begin{cases}2D & \text{(complex-valued or imag. relation embeddings)} \\ D & \text{(real-valued)}\end{cases}
\label{eq:dim}
\end{equation}
The Transformer hidden dimension is $d_{\text{model}} = \max(d_{\mathcal{V}}, d_{\mathcal{R}})$.
Two learnable linear projections align token dimensions:
\begin{equation}
\operatorname{Proj}_{\mathcal{V}}: \mathbb{R}^{d_{\mathcal{V}}} \to \mathbb{R}^{d_{\text{model}}},\qquad
\operatorname{Proj}_{\mathcal{R}}: \mathbb{R}^{d_{\mathcal{R}}} \to \mathbb{R}^{d_{\text{model}}}
\label{eq:projections}
\end{equation}
Both are learned during PAMR training while the KGE model remains frozen;
$\operatorname{Proj}_{\mathcal{V}}$ reduces to the identity when $d_{\mathcal{V}} = d_{\text{model}}$.

\vspace{-1em}
\subsection{Path Extraction}
\vspace{-1em}
\label{sec:path-extraction}

To provide the Transformer with explicit, interpretable connectivity evidence rather than
relying solely on the implicit proximity captured by KGE embeddings, we extract shortest
paths between query pairs via \textbf{bidirectional breadth-first search} (BFS).  The
core idea is to launch two
symmetric search waves, one forward from $\mathbf{s}$ and one backward from $\mathbf{o}$, and
expand them layer by layer until their frontiers intersect.  At each depth $k$, let
$\mathcal{F}_s^{k}$ and $\mathcal{F}_o^{k}$ denote the forward and backward frontier vertex
sets.  The frontiers expand according to:
\begin{align}
\mathcal{F}_s^{k} &= \{\, \mathbf{v} \in \mathcal{V} \mid \exists\, \mathbf{w} \in \mathcal{F}_s^{k-1},\;
   (\mathbf{w}, r, \mathbf{v}) \in \mathcal{E} \,\} \label{eq:fw-frontier} \\
\mathcal{F}_o^{k} &= \{\, \mathbf{v} \in \mathcal{V} \mid \exists\, \mathbf{w} \in \mathcal{F}_o^{k-1},\;
   (\mathbf{v}, r, \mathbf{w}) \in \mathcal{E} \,\} \label{eq:bw-frontier}
\end{align}
When the intersection $\mathcal{I}_k = \mathcal{F}_s^{k} \cap \mathcal{F}_o^{k}$ is non-empty for
the first time, the meeting vertex $m$ yields a path $p = \text{pred}_s(m) \parallel \text{rev}(\text{pred}_o(m))$
that is guaranteed to be minimal in hop count by construction.


\vspace{-1em}
\paragraph{Decaying fanout heuristic.}
To remain scalable on large, dense knowledge graphs, we employ a \emph{decaying fanout}
mechanism~\cite{nam2025sequential} that prunes the search space by limiting the number of
neighbors expanded at each depth.  With a base fanout $F$ and decay rate $\lambda \in (0, 1]$,
the effective fanout at depth $k$ from each side is $\lfloor F \cdot \lambda^{k} \rfloor$.
For each node $\mathbf{u}$ at depth $k$, only the $\lfloor F \cdot \lambda^{k} \rfloor$
highest-frequency incident edges are expanded, favouring well-connected paths.  The full
algorithm is detailed in Algorithm~\ref{alg:bfs}.

\begin{algorithm}[htbp]
\caption{Bidirectional BFS with Decaying Fanout}
\label{alg:bfs}
\KwIn{Subject $s$, candidate object $\hat{o}$, graph $\mathcal{G}$, 
      base fanout $F$, decay rate $\lambda \in (0,1]$, max hops $H$}
\KwOut{Shortest path from $s$ to $\hat{o}$ within $H$ hops, or $\emptyset$ if unreachable}

$\mathcal{Q}_s \gets \{s\}$; $\mathcal{Q}_o \gets \{\hat{o}\}$\;
$\mathtt{pred}_s[s] \gets \bot$; $\mathtt{pred}_o[\hat{o}] \gets \bot$\;
$\mathtt{visited}_s \gets \{s\}$; $\mathtt{visited}_o \gets \{\hat{o}\}$\;
$\mathtt{depth}_s[s] \gets 0$; $\mathtt{depth}_o[\hat{o}] \gets 0$\;
\While{$\mathcal{Q}_s \neq \emptyset$ \textbf{and} $\mathcal{Q}_o \neq \emptyset$}{
    \uIf{$|\mathcal{Q}_s| \leq |\mathcal{Q}_o|$}{
        $m \gets$ \texttt{expandFrontier}($\mathcal{Q}_s$, $\mathtt{visited}_s$, $\mathtt{visited}_o$, 
            $\mathtt{pred}_s$, $\mathtt{depth}_s$, $F$, $\lambda$, \textit{forward})\;
    }
    \Else{
        $m \gets$ \texttt{expandFrontier}($\mathcal{Q}_o$, $\mathtt{visited}_o$, $\mathtt{visited}_s$, 
            $\mathtt{pred}_o$, $\mathtt{depth}_o$, $F$, $\lambda$, \textit{backward})\;
    }
    \If{$m \neq \bot$}{
        \Return $\texttt{reconstructPath}(m, \mathtt{pred}_s, \mathtt{pred}_o)$\;
    }
    \If{$\mathtt{depth}_s \geq H$ \textbf{or} $\mathtt{depth}_o \geq H$}{
        \Return $\emptyset$\;
    }
}
\Return $\emptyset$\;

\SetKwProg{Fn}{function}{:}{}
\SetKwFunction{FExpand}{expandFrontier}
\Fn{\FExpand{$\mathcal{Q}$, $\mathtt{visitedSelf}$, $\mathtt{visitedOpp}$, 
    $\mathtt{pred}$, $\mathtt{depth}$, $F$, $\lambda$, $\mathit{direction}$}}{
    \ForEach{$\mathbf{u}$ in current depth-level of $\mathcal{Q}$}{
        $k \gets \mathtt{depth}[\mathbf{u}]$\;
        $\mathit{fanout} \gets \lfloor F \cdot \lambda^{k} \rfloor$\;
        $\mathit{edges} \gets$ top $\mathit{fanout}$ most frequent incident edges to $\mathbf{u}$ 
                  (adjusted for $\mathit{direction}$)\;
        \ForEach{$(\mathbf{w}, r, \mathbf{v})$ in $\mathit{edges}$}{
            \If{$\mathbf{v} \notin \mathtt{visitedSelf}$}{
                $\mathtt{visitedSelf} \gets \mathtt{visitedSelf} \cup \{\mathbf{v}\}$\;
                $\mathcal{Q}.\text{enqueue}(\mathbf{v})$\;
                $\mathtt{pred}[\mathbf{v}] \gets (\mathbf{u}, r, \mathit{direction})$\;
                $\mathtt{depth}[\mathbf{v}] \gets k + 1$\;
            }
            \If{$\mathbf{v} \in \mathtt{visitedOpp}$}{
                \Return $\mathbf{v}$\;
            }
        }
    }
    \Return $\bot$\;
}
\end{algorithm}

\vspace{-1em}
\paragraph{Complexity analysis.}
The search with decaying fanout visits at most $\lfloor F \cdot \lambda^{k} \rfloor$ neighbors
per node at depth $k$.  The total number of nodes visited by one search wave is bounded by
the geometric series $\sum_{k=0}^{K/2} F \cdot \lambda^{k} < \frac{F}{1 - \lambda}$,
for $\lambda < 1$, yielding space complexity $O\bigl(\frac{F}{1-\lambda}\bigr)$ and
time complexity $O\bigl(\bar{d} \cdot \frac{F}{1-\lambda}\bigr)$, with $\bar{d}$ denoting the
average degree and $K$ the shortest path length.  This is exponentially more efficient than
the unconstrained bidirectional BFS complexity of $O(\bar{d}^{\,K/2 + 1})$, at the cost of
relaxing the optimality guarantee, meaning a path may be missed if pruned by the fanout heuristic.

\vspace{-1em}
\paragraph{PAMR vs.\ PAR graph configurations.}
We define two graph configurations that produce qualitatively different path spaces:
\begin{description}[nosep]
    \item[$\mathcal{G}_{\text{full}}$ (PAMR).]  The complete graph comprising all entity types.
    Paths may bridge across artworks, artists, and attributes, exploiting cross-type
    relationships for reasoning.
    \item[$\mathcal{G}_{\text{artist}}$ (PAR).]  The artist-only subgraph restricted to
    entities of type~2 and edges where both endpoints are artists.  Paths stay within the
    artist community and leverage only relations between artists.
\end{description}
Paths with $\ell \leq 0$ are excluded from training and evaluation as they carry no
structural reasoning signal, which include unreachable pairs and trivial self loops.

\vspace{-1em}
\subsection{PAMR Architecture}
\vspace{-1em}
\label{sec:pamr-architecture}

\vspace{-1em}
\subsubsection{Token Interleaving}
\vspace{-1em}
\label{sec:token-interleaving}

To encode path structure as a flat sequence, entity and relation tokens are interleaved
with a classification token and the query relation prepended.
Let $\mathbf{e}^{\text{raw}}_i \in \mathbb{R}^{d_{\mathcal{V}}}$ be the raw KGE embedding
of entity $i$ and $\mathbf{r}^{\text{raw}}_j \in \mathbb{R}^{d_{\mathcal{R}}}$ the raw
relation embedding of relation $j$.  Projected token embeddings are:
\begin{equation}
\mathbf{v}_i = \operatorname{Proj}_{\mathcal{V}}(\mathbf{e}^{\text{raw}}_i),\qquad
\mathbf{r}_j = \operatorname{Proj}_{\mathcal{R}}(\mathbf{r}^{\text{raw}}_j)
\label{eq:token-proj}
\end{equation}
both lying in $\mathbb{R}^{d_{\text{model}}}$.

A path with $\ell$ hops comprises $\ell+1$ entity tokens and $\ell$ edge tokens.
These are interleaved into a token sequence of length $L = 2\ell + 1$:
\begin{equation}
\mathbf{T} = \bigl[\mathbf{c};\; \mathbf{r}_q;\;
                     \mathbf{v}_{\mathbf{s}};\; \mathbf{r}_{1};\; \mathbf{v}_{2};\; \mathbf{r}_{2};\;
                     \mathbf{v}_{3};\; \ldots;\; \mathbf{v}_{\mathbf{o}}\bigr]
\in \mathbb{R}^{L \times d_{\text{model}}}
\label{eq:token-seq}
\end{equation}
Here $\mathbf{c} \in \mathbb{R}^{d_{\text{model}}}$ is a learnable classification token
at position 0, initialised from $\mathcal{N}(0, 0.02^2)$.  $\mathbf{r}_q$ is the query
relation representation at position 1, while entity tokens $\mathbf{v}_{\mathbf{s}}, \mathbf{v}_{2}, \ldots, \mathbf{v}_{\mathbf{o}}$
occupy positions $2 + 2i$ for $i = 0, \ldots, \ell$, and edge tokens $\mathbf{r}_{1}, \ldots, \mathbf{r}_{\ell}$
occupy positions $3 + 2i$ for $i = 0, \ldots, \ell - 1$.

Entity positions receive additive \textbf{node type embeddings}
$\boldsymbol{\tau}_i \in \mathbb{R}^{d_{\text{model}}}$ from a learned embedding table
with one entry per node type plus a padding zero.
All positions receive \textbf{absolute positional embeddings}
$\mathbf{p}_i = \operatorname{PE}(i)$ via a learned embedding
$\operatorname{PE}: \{0, \ldots, L-1\} \to \mathbb{R}^{d_{\text{model}}}$.
The final token input to the Transformer is:
\begin{equation}
\mathbf{T}^{\text{in}}_i = \mathbf{T}_i \;+\; \mathbf{p}_i \;+\;
    \boldsymbol{\tau}_{\tau(i)} \cdot \mathds{1}_{\mathcal{V}}(i)
\label{eq:token-in}
\end{equation}
Here $\mathds{1}_{\mathcal{V}}(i) = 1$ iff position $i$ is an entity position.

\vspace{-1em}
\subsubsection{Entity and Relation Representations}
\vspace{-1em}
\label{sec:entity-relation-repr}

Entity tokens embed only the fused real part (Eq.~\ref{eq:fused-emb}):
$\mathbf{e}^{\text{raw}}_i = \mathbf{e}^{\text{fused}}_{i,\text{re}}$.
Relation tokens concatenate real and imaginary components according to Eq.~\ref{eq:rel-raw}.

\vspace{-1em}
\subsubsection{Transformer Encoder}
\vspace{-1em}
\label{sec:transformer-encoder}

The token sequence is contextualised through a stack of $L_{\text{tr}}$
Transformer encoder layers~\cite{vaswani2017attention} with $h$ attention heads each:
\begin{equation}
\mathbf{H}^{(0)} = \mathbf{T}^{\text{in}},\qquad
\mathbf{H}^{(l)} = \operatorname{EncLayer}^{(l)}\bigl(\mathbf{H}^{(l-1)},\; \mathbf{M}\bigr)
\quad (l = 1, \ldots, L_{\text{tr}})
\label{eq:transformer}
\end{equation}
Here $\mathbf{M} \in \{0, 1\}^{L}$ is a padding mask ($\mathbf{M}_i = 1$ indicates an
invalid position) for paths shorter than the batch maximum $N_{\max}(\mathcal{B})$.
Each encoder layer applies:
\begin{align}
    \tilde{\mathbf{H}} &= \operatorname{LayerNorm}\bigl(
        \mathbf{H}^{(l-1)} + \operatorname{MHA}(\mathbf{H}^{(l-1)}; \mathbf{M})\bigr)
        \label{eq:attn}\\
    \mathbf{H}^{(l)} &= \operatorname{LayerNorm}\bigl(
        \tilde{\mathbf{H}} + \operatorname{FFN}(\tilde{\mathbf{H}})\bigr)
        \label{eq:ffn}
\end{align}
with $\operatorname{FFN}(\mathbf{x}) = \mathbf{W}^{(2)} \operatorname{GELU}(\mathbf{W}^{(1)} \mathbf{x} + \mathbf{b}^{(1)}) + \mathbf{b}^{(2)}$
and inner dimension $d_{\text{ff}} = 4 \times d_{\text{model}}$.  GELU~\cite{hendrycks2016gaussian} activation
and Post-LN layer normalisation are used throughout.

\vspace{-1em}
\subsubsection{Readout and Scoring}
\vspace{-1em}
\label{sec:readout}

To obtain a fixed-dimensional path representation, the contextualised token
sequence is aggregated via one of two readout strategies from
$\mathbf{H}^{(L_{\text{tr}})} \in \mathbb{R}^{L \times d_{\text{model}}}$:

\begin{description}[nosep]
    \item[CLS.] $\mathbf{z} = \mathbf{H}^{(L_{\text{tr}})}[0, :] \in \mathbb{R}^{d_{\text{model}}}$, the classification token at position 0.
    \item[Mean.] $\mathbf{z} = \frac{\sum_{i=0}^{L-1} \mathbf{H}^{(L_{\text{tr}})}[i, :] \cdot (1 - \mathbf{M}[i])}{\sum_i (1 - \mathbf{M}[i])}$, averaging over valid tokens.
\end{description}

A two-layer MLP maps $\mathbf{z}$ to a scalar score $s \in \mathbb{R}$:
\begin{equation}
    \mathbf{h} = \operatorname{GELU}\bigl(\mathbf{W}_1 \mathbf{z} + \mathbf{b}_1\bigr),
    \qquad
    s = \mathbf{w}_2^{\top} \mathbf{h} + b_2
\label{eq:readout}
\end{equation}
with dropout $p = 0.1$ applied after GELU.  All linear layers use Xavier
uniform initialisation and zero bias.

\vspace{-1em}
\subsection{Training}
\vspace{-1em}
\label{sec:training}

To train a path scorer with stable supervision, each query must contribute equal gradient
signal regardless of its number of negative samples.  We construct a paired sampling
dataset that guarantees both balance and fairness.

\vspace{-1em}
\subsubsection{Paired Sampling}
\vspace{-1em}
\label{sec:paired-sampling}

For each query $q$, the query generation process produces one \textit{positive} path
$p^+_q$, which connects $\mathbf{s}$ to the true object $\mathbf{o}$, and $K$ \textit{negative}
paths $\mathcal{S}^-_q = \{p^-_{q,1}, \ldots, p^-_{q,K}\}$, which connect $\mathbf{s}$ to
sampled false objects.  We construct a paired dataset $\mathcal{D}_{\text{pair}}$
that iterates over every negative exactly once per epoch, each paired with its query's
unique positive.  The construction procedure is detailed in Algorithm~\ref{alg:paired}.

\begin{algorithm}[htbp]
\caption{Paired-Sampling Dataset Construction}
\label{alg:paired}
\KwIn{Path samples grouped by query identifier $q$: $\mathcal{S}_q^+ = \{p^+_q\}$,
      $\mathcal{S}_q^- = \{p^-_{q,1}, \ldots, p^-_{q,K}\}$}
\KwOut{Paired index $\mathcal{I}$; $|\mathcal{I}| = \sum_q |\mathcal{S}_q^-|$}

$\mathcal{I} \gets \emptyset$\;
\ForEach{query $q$}{
    \ForEach{$p^- \in \mathcal{S}_q^-$}{
        $\mathcal{I}.\text{append}(\,(p^+_q,\; p^-)\,)$\;
    }
}
\Return $\mathcal{I}$\;
\end{algorithm}

Collation interleaves positive and negative pairs into flat sequences:
$\operatorname{batch}(B) = [p^+_0, p^-_0, p^+_1, p^-_1, \ldots, p^+_{B-1}, p^-_{B-1}]$.
Here $B$ is the query batch size.  Each training batch contains exactly $B$ positive and
$B$ negative samples, producing a guaranteed 50/50 composition regardless of the underlying
positive to negative ratio.  Variable length paths are zero-padded to the batch maximum
$N_{\max}(\mathcal{B}) = \max_{p \in \mathcal{B}} (\ell_p + 1)$ entity tokens, with
padding masks $\mathbf{M}$ as defined in Section~\ref{sec:token-interleaving}.

\vspace{-1em}
\subsubsection{Loss Functions}
\vspace{-1em}
\label{sec:loss-functions}

The architecture supports three loss modes to accommodate different dataset
characteristics.  Let $\mathcal{B} = \{(s_1, y_1), \ldots, (s_{2B}, y_{2B})\}$
denote a training batch with $y_i \in \{0, 1\}$.

\vspace{-1em}
\paragraph{Binary Cross Entropy.}  A pointwise loss treating each path independently:
\begin{equation}
\mathcal{L}_{\text{BCE}} = -\frac{1}{2B} \sum_{(s_i, y_i) \in \mathcal{B}}
    \bigl[\, y_i \log \sigma(s_i) + (1 - y_i) \log(1 - \sigma(s_i)) \,\bigr]
\label{eq:loss-bce}
\end{equation}
Here $\sigma(x) = 1/(1 + e^{-x})$ denotes the sigmoid function.  With perfectly balanced batches,
the loss converges to $\ln 2 \approx 0.693$ when the model predicts $\sigma(s_i) = 0.5$
for all inputs, which is a degenerate state yielding $\mathrm{MRR}=1.0$ artifactually.

\vspace{-1em}
\paragraph{Link Prediction (softplus).}  The standard KGE loss~\cite{trouillon2016complex}:
\begin{equation}
\mathcal{L}_{\text{LP}} = \frac{1}{2B} \sum_{(s_i, y_i) \in \mathcal{B}}
    \log\bigl(1 + \exp(-\hat{y}_i \cdot s_i)\bigr),\quad
    \hat{y}_i = 2y_i - 1 \in \{-1, +1\}
\label{eq:loss-lp}
\end{equation}

\vspace{-1em}
\paragraph{Ranking (margin).}  In contrast to the pointwise losses above, this loss
groups samples by query identifier and enforces a margin
between the positive score and all negative scores for that query:
\begin{equation}
\mathcal{L}_{\text{rank}} = \frac{1}{|\mathcal{Q}_{\mathcal{B}}|} \sum_{q \in \mathcal{Q}_{\mathcal{B}}}\;
    \frac{1}{|S_q^+||S_q^-|} \sum_{s^+ \in S_q^+} \sum_{s^- \in S_q^-}
    \max\bigl(0,\, m - s^+ + s^-\bigr)
\label{eq:loss-rank}
\end{equation}
Here $\mathcal{Q}_{\mathcal{B}}$ denotes the set of unique query identifiers in $\mathcal{B}$, and
$m = 1.0$ is the margin.  This loss directly optimises the ranking objective and prevents
score collapse.

\vspace{-1em}
\subsubsection{Optimization}
\vspace{-1em}
\label{sec:optimization}

PAMR is trained with AdamW~\cite{loshchilov2019decoupled}:
$\eta = 0.001$, $\lambda = 0.01$, $\beta_1 = 0.9$, $\beta_2 = 0.999$.
Only PAMR parameters are updated, while the frozen VL-KGE model remains untouched.

\vspace{-1em}
\subsection{Evaluation}
\vspace{-1em}

For each query $(\mathbf{s}, r, ?)$ the model assigns a scalar score to every
candidate object and ranks them in descending order.  Let $r_q$ denote the
position of the true object in this ranking for query $q$, out of
$|\mathcal{Q}|$ total queries.  Standard ranking metrics are computed:
\begin{equation}
    \mathrm{MRR} = \frac{1}{|\mathcal{Q}|}\sum_{q} \frac{1}{r_q},
    \qquad
    \mathrm{Hits@}K = \frac{1}{|\mathcal{Q}|}\sum_{q} \mathds{1}[r_q \leq K].
\end{equation}

\vspace{-1em}
\paragraph{Per-relation analysis.}
To assess performance independently for each relation type, all metrics are
disaggregated by $r \in \mathcal{R}$.  Let $\mathcal{Q}_r = \{q : \operatorname{rel}(q) = r\}$
be the subset of queries whose relation is $r$.  Then
$\mathrm{MRR}_r = \frac{1}{|\mathcal{Q}_r|}\sum_{q \in \mathcal{Q}_r} 1/r_q$
and $\mathrm{Hits@}K_r = \frac{1}{|\mathcal{Q}_r|}\sum_{q \in \mathcal{Q}_r} \mathds{1}[r_q \leq K]$.

\vspace{-1em}
\subsection{Baselines}
\vspace{-1em}
\label{sec:baselines}

\vspace{-1em}
\paragraph{VL-KGE models.}
To establish a performance floor and isolate the marginal contribution of explicit path
reasoning, we compare PAMR against four pre-trained KGE baselines using the same frozen
embeddings.  These models score candidate triples
$(\mathbf{s}, r, \mathbf{o})$ via learned embeddings without path reasoning.
All models share identical training configuration: 20 epochs, $\eta = 0.1$, batch size
$4096$, $K = 1$ negative per positive triple, and multimodal fusion (average of structural,
visual, and textual features).

\begin{table}[t]
\centering
\caption{Baseline VL-KGE models.  $D = 768$. Parameter counts exclude frozen feature extractors.}
\label{tab:kge-models}
\begin{tabular}{lcc}
\toprule
\textbf{Model} & \textbf{Scoring $\phi(\mathbf{s}, r, \mathbf{o})$} & \textbf{Params} \\
\midrule
ComplEx & $\operatorname{Re}\langle \mathbf{s}, \mathbf{r}, \overline{\mathbf{o}} \rangle$
        & 350.8M \\
TransE  & $-\|\mathbf{s} + \mathbf{r} - \mathbf{o}\|_1$
        & 172.1M \\
DistMult & $\sum_i \mathbf{s}_i \mathbf{r}_i \mathbf{o}_i$
        & 172.1M \\
RotatE  & $-\|\mathbf{s} \circ \mathbf{r} - \mathbf{o}\|$
        & 350.8M \\
\bottomrule
\end{tabular}
\end{table}

All models are evaluated on the same held-out test set using identical evaluation
metrics, MRR and Hits@$K$, enabling fair comparison with PAMR.  The KGE baselines
represent an upper bound without relying on path structure: they leverage the full
knowledge graph structure through learned embeddings but cannot exploit explicit
path topology between query pairs.  PAMR must learn to attend over the paths to
match or exceed this signal.
